\chapter*{序言：结构、活动与持续生成}
\addcontentsline{toc}{chapter}{序言：结构、活动与持续生成}
\label{chp:prelude_geometry}
我们理解智能，不能只看某一时刻的输出，也不能只看保存下来的知识。当前活动决定系统此刻能够处理什么，历史影响它如何解释输入，结构规定可达的路径，目标决定哪些路径值得探索。智能的连续性来自这些过程的共同演化。

河道与流水可以提供一个有限的直觉：结构约束活动，活动在一定条件下改变结构。我们借助这个意象说明相互作用，但不会由它推断思维具有水的物性。计算模型必须进一步给出状态变量、更新规则和边界条件。

同样，自我可以被讨论为系统维护的连续约束与自我模型。其是否对应主观体验，需要认知科学与哲学层面的独立论证。我们不从形式相似性推导意识，也不从模型中的几何轨迹推导个体生命之外的持续存在。

本书由此出发：把富有启发性的意象转化为能够检验的机制，把宏观解释连接到局部更新，把持续生成连接到有限资源下的具体运行。
